\chapter{EM for the transition law from noised sequences}
\label{ch:em-through-noise}

This chapter assembles the paper's estimator from the pieces of the previous four. The
data are $\nseq$ noised sequences, each observed once at its own noise level; the
unknown is the transition kernel $\kernelth$ of the clean chain; the objective is the
marginal likelihood; the algorithm is expectation--maximisation with the exact E-step
of \cref{ch:factor-graphs-bp} on the grid of \cref{ch:grid-quadrature}. Everything the
M-step needs is one $\ngrid\times\ngrid$ matrix, the expected transition statistic
$\Xis$, and Fisher's identity turns the same matrix into the exact gradient of the
marginal likelihood. The chapter derives Fisher's identity in full, defines $\Xis$ and
its properties, derives the Gaussian M-step (the belief-weighted Yule--Walker
equations), describes the mixture-innovation kernel and its conditional-maximisation
sweeps, explains the checkpoint ladder, the validation selection and the
\code{at\_cap} flag, the clean-data limit, and how E06, E06c and E02a use all of it. It
ends with the information-theoretic reason for the convergence asymmetry the paper
reports.

\section{The marginal likelihood of noised sequences}

The paper's Section 4 drops the assumption that $\kernel$ is known and keeps
everything else: $a_1\sim\normal01$, $a_i=\corr a_{i-1}+\varepsilon_i$ with a
parameterised transition kernel $\kernelth(a'\mid a)$, and the channel
$\vx=e^{-t}\va+\sqrt{\Dt}\vz$. The observations are $\nseq$ sequences $\vx^\mu$, the
$\mu$-th at level $t_\mu$. Under the \term{one-view protocol} each clean sequence is
noised once, at one level, and never seen clean; in the headline experiments the
sequences are split evenly across the twelve levels of the schedule, and in the
learning experiment E06 all of them sit at one level. The \term{marginal
log-likelihood} is
\begin{equation}
  \Lik(\params)=\sum_{\mu=1}^{\nseq}\log\Pt[t_\mu,\params](\vx^\mu),\qquad
  \Pt[t,\params](\vx)=\int\mu(a_1)\prod_{i=2}^{\nsites}\kernelth(a_i\mid a_{i-1})
  \prod_{i=1}^{\nsites}\like(x_i\mid a_i)\dd\va,
  \label{eq:en-marginal}
\end{equation}
where the unary factors $\like$ depend on $t_\mu$ but not on $\params$, and only the
transition factors carry $\params$. Each term of \cref{eq:en-marginal} is the evidence
of \cref{ch:factor-graphs-bp}, returned by the forward sweep on the grid. Two structural
facts drive everything below: the parameter enters only through the $\nsites-1$
transition factors of each sequence, and the noise level enters only through the unary
factors, so sequences at different levels contribute additively and identically to
every statistic (\code{e\_step\_multi} in \code{src/markov\_score/em.py} simply adds
them).

\section{Fisher's identity}

\subsection{Derivation}

Let $p_\params(x,z)$ be a joint density of data and latent variables, with
$p_\params(x)=\int p_\params(x,z)\dd z$, and assume differentiation under the integral
sign is permitted (the integrand and its $\params$-derivative are dominated by an
integrable function, which holds for the kernels used here on the truncated grid, where
the integral is a finite sum). Then
\begin{equation}
\begin{aligned}
  \nabla_\params\log p_\params(x)&=\frac{\nabla_\params p_\params(x)}{p_\params(x)}
  =\frac{1}{p_\params(x)}\int\nabla_\params p_\params(x,z)\dd z\\
  &=\int\frac{p_\params(x,z)}{p_\params(x)}\,\nabla_\params\log p_\params(x,z)\dd z
  =\E_{p_\params(z\mid x)}\bigl[\nabla_\params\log p_\params(x,z)\bigr],
\end{aligned}
  \label{eq:en-fisher-general}
\end{equation}
using $\nabla p=p\nabla\log p$ in the third step and $p_\params(x,z)/p_\params(x)=p_\params(z\mid x)$
in the last. This is \term{Fisher's identity}: the gradient of the marginal
log-likelihood equals the posterior expectation of the gradient of the complete-data
log-likelihood. It was met in \cref{ch:maximum-likelihood-em} as the statement that
EM's fixed points are stationary points; here it is a computational device.

Apply it to one sequence of \cref{eq:en-marginal}. The complete-data log-likelihood is
$\log\mu(a_1)+\sum_{i\ge2}\log\kernelth(a_i\mid a_{i-1})+\sum_i\log\like(x_i\mid a_i)$, and
only the middle sum depends on $\params$, so
\begin{equation}
\begin{aligned}
  \nabla_\params\log\Pt[t,\params](\vx)
  &=\sum_{i=2}^{\nsites}\E\bigl[\nabla_\params\log\kernelth(a_i\mid a_{i-1})\,\big|\,\vx\bigr]\\
  &=\sum_{i=2}^{\nsites}\int p_t(a_{i-1},a_i\mid\vx)\,\nabla_\params\log\kernelth(a_i\mid a_{i-1})\dd a_{i-1}\dd a_i.
\end{aligned}
  \label{eq:en-fisher-chain}
\end{equation}
Each term needs only the \emph{pairwise} posterior marginal on one edge, which the
two sweeps return exactly (equation \emph{pair} of \cref{ch:factor-graphs-bp}). No
derivative of the recursion is taken: the sweeps supply an expectation, and the
gradient is an integral against it. This is the sense of the paper's ``no
differentiation through the recursion is needed''.

\subsection{On the grid: the expected transition statistic}

Replace the integrals by the discretised model of \cref{ch:grid-quadrature}. The
pairwise posterior on edge $(i-1,i)$ of sequence $\mu$ becomes a matrix of
probabilities $P(a_i=\gridpt{k},a_{i-1}=\gridpt{l}\mid\vx^\mu)$ over grid pairs, and
$\nabla_\params\log\kernelth$ becomes the array $\nabla_\params\log\kernelth(\gridpt{k}\mid\gridpt{l})$.
Summing over edges and sequences,
\begin{equation}
\begin{aligned}
  \nabla_\params\Lik(\params)&=\sum_{k,l}\Xis_{kl}(\params)\,\nabla_\params\log\kernelth(\gridpt{k}\mid\gridpt{l})
  =\bigl\langle\Xis(\params),\nabla_\params\log\kernelth\bigr\rangle,\\
  \Xis_{kl}&=\sum_\mu\sum_{i\ge2}P(a_i=\gridpt{k},a_{i-1}=\gridpt{l}\mid\vx^\mu),
\end{aligned}
  \label{eq:en-fisher-grid}
\end{equation}
which is the paper's equation \emph{fisher}. The matrix $\Xis$ is the
\term{expected transition statistic}: entry $(k,l)$ is the expected posterior mass on
the transition $\gridpt{l}\to\gridpt{k}$ (column the parent, row the child, the same
orientation as the transition matrix), summed over every edge of every sequence. It is
the continuous-state analogue of the expected transition counts of Baum--Welch. Its
computation from the sweeps is a rank-one accumulation: with
$f_i=w\circ\lmsg_i\circ\ell_i$ and $g_i=w\circ\rmsg_i\circ\ell_i$ (messages with the local
factor and the quadrature weight folded in), the unnormalised pairwise belief on edge
$(i,i+1)$ is $f_i(l)\,\kernel[k,l]\,g_{i+1}(k)$ with normaliser $Z_i=g_{i+1}^{\!\top}\kernel f_i$, so
\begin{equation}
  \Xis=C\circ\kernel,\qquad C=\sum_{\mu}\sum_{i}\frac{g^\mu_{i+1}\,(f^\mu_i)^{\!\top}}{Z^\mu_i},
  \label{eq:en-xi-accum}
\end{equation}
with $\circ$ the elementwise product. The code (\code{\_e\_step\_chunk} in
\code{src/markov\_score/em.py}) forms $C$ as one matrix product of the stacked
$g$ and $f$ vectors over all edges of a chunk, then multiplies by $\kernel$.

Three properties of $\Xis$ are used constantly.
\begin{enumerate}
\item \emph{Mass.} $\sum_{kl}\Xis_{kl}=\nseq(\nsites-1)$, the number of transitions,
because each pairwise posterior sums to one. This is checked exactly in
\code{tests/test\_em\_bp.py} and is what lets the clean-data histogram of
\cref{sec:en-clean} be a drop-in replacement.
\item \emph{Additivity.} Sequences at different levels contribute additively; one
$\Xis$ summarises a dataset spread over the whole schedule.
\item \emph{What it is not.} $\Xis$ is a complete summary for the current M-step (the
maximisation never revisits the data), but it is not a sufficient statistic of the
observations: it is computed at the current $\params$ and depends on it, on the grid,
on the fixed initial law and on the homogeneity of the transition. The code's
docstring insists on this distinction, and so should the reader.
\end{enumerate}

\subsection{The same matrix is the whole E-step}

With $\Xis$ in hand, the Q-function of \cref{ch:maximum-likelihood-em} for the
discretised model is
\begin{equation}
  \Qfun(\params\mid\params')=\bigl\langle\Xis(\params'),\log\kernelth\bigr\rangle+\text{(site-1 term)}+\text{const},
  \label{eq:en-Q}
\end{equation}
a finite bilinear form; the site-1 term involves the initial law, which the paper
holds fixed at $\normal01$, and the constant collects the likelihood terms and the
posterior entropy, none of which depend on $\params$. Its gradient at $\params=\params'$
is \cref{eq:en-fisher-grid}: the Q-function's gradient and the marginal likelihood's
gradient coincide where the E-step was taken, which is Fisher's identity once more.
The E-step is therefore ``one forward--backward pass'' plus the accumulation
\cref{eq:en-xi-accum}, at cost $O(\nseq\nsites\ngrid^2)$ independent of the number of
parameters, and the M-step touches only $\Xis$, at cost $O(P\ngrid^2)$ for $P$
parameters.

\subsection{Verification}

The paper states that the identity is verified against finite differences to
$10^{-9}$. Three tests do the work, and it matters which checks what.
\code{test\_fisher\_identity\_matches\_finite\_differences} in
\code{tests/test\_em\_bp.py} compares $\langle\Xis,\nabla\log\kernelth\rangle$ with a
central difference of the grid evidence for the Gaussian kernel: a real check, but not
an independent one, since the evidence and $\Xis$ come from the same recursion and the
same discretised kernel, so a shared error (a transposed index, a misplaced factor)
could satisfy both sides. \code{tests/test\_fisher\_bruteforce.py} removes that
coupling by enumerating every latent path of a five-state, four-site chain from the
model definition, with no recursion, and checking the evidence, $\Xis$ \emph{including
its orientation} (the matrix must match and must not match its transpose, on a grid
asymmetric enough that the two differ), and the gradient against a finite difference of
the \emph{enumerated} evidence. The Laplace kernel is included there deliberately: its
$\corr$-derivative carries $\operatorname{sign}(a'-\corr a)$, a discontinuity along the
diagonal, and trapezoidal quadrature of a discontinuous integrand is first order
(\cref{ch:grid-quadrature}); the test in \code{tests/test\_em\_bp.py} named for the
Laplace scale gradient pins the fact that the smooth $b$-derivative matches a finite
difference to high accuracy while the $\corr$-derivative is off by a grid-dependent
percentage. That is a discretisation artefact of differentiating a non-smooth density
on a mesh, not a failure of the identity; the Laplace M-step avoids it by never using
the gradient (it is a weighted median). Finally,
\code{tests/test\_gradient\_trace\_alignment.py} guards a different invariant of the
gradient-ascent route: the returned parameter is the one at which the last recorded
likelihood was evaluated, and the monotonicity record covers the final update.

\section{The M-step for the Gaussian kernel}

Take $\kernelth(a'\mid a)=\normal{\corr a}{\ivar}$, $\params=(\corr,\ivar)$. The edge part
of \cref{eq:en-Q} is
\begin{equation}
  \Qfun(\corr,\ivar)=\sum_{k,l}\Xis_{kl}\Bigl[-\frac{(\gridpt{k}-\corr\gridpt{l})^2}{2\ivar}-\frac12\log(2\pi\ivar)\Bigr]
  =-\frac{S_{11}-2\corr S_{01}+\corr^2S_{00}}{2\ivar}-\frac{W}{2}\log(2\pi\ivar),
  \label{eq:en-Qgauss}
\end{equation}
where $W=\sum_{kl}\Xis_{kl}$ and the three belief-weighted moments are
\[
  S_{00}=\sum_{kl}\Xis_{kl}\,\gridpt{l}^2,\qquad
  S_{01}=\sum_{kl}\Xis_{kl}\,\gridpt{k}\gridpt{l},\qquad
  S_{11}=\sum_{kl}\Xis_{kl}\,\gridpt{k}^2,
\]
the second moment of the parent, the cross moment, and the second moment of the child.
The subscripts index the two variables of a transition ($0$ for the parent
$\gridpt{l}$, $1$ for the child $\gridpt{k}$) and say which pair is multiplied; they are
not powers. The paper's Section 4 writes the same three quantities as
$S_{ab}=\sum_{kl}\Xis_{kl}\gridpt{k}^a\gridpt{l}^b$, which read literally as powers would
make $S_{00}=W$ and $S_{01}=\sum\Xis_{kl}\gridpt{l}$ and would not give the formulas that
follow; the index reading above is the one the formulas and the code use.
Setting the derivatives to zero,
\begin{equation}
\begin{aligned}
  \frac{\partial\Qfun}{\partial\corr}&=\frac{S_{01}-\corr S_{00}}{\ivar}=0
  &&\Rightarrow\quad \corr=\frac{S_{01}}{S_{00}},\\
  \frac{\partial\Qfun}{\partial\ivar}&=\frac{S_{11}-2\corr S_{01}+\corr^2S_{00}}{2\ivar^2}-\frac{W}{2\ivar}=0
  &&\Rightarrow\quad \ivar=\frac{S_{11}-2\corr S_{01}+\corr^2S_{00}}{W}=\frac{S_{11}-\corr S_{01}}{W},
\end{aligned}
  \label{eq:en-yw}
\end{equation}
the last equality because $\corr^2S_{00}=\corr S_{01}$ at the optimum. These are the
\term{belief-weighted Yule--Walker equations}: the least-squares regression of the
child on the parent through the origin, with the pairwise posterior mass as weights,
and the residual variance as the innovation variance. The Hessian of
\cref{eq:en-Qgauss} is negative definite at this point (it is a weighted least-squares
problem in $\corr$ and a concave profile in $\ivar$), so this is the maximiser and the
M-step is exact. In the clean-data limit, where $\Xis$ is a histogram of observed
transitions, \cref{eq:en-yw} reduces to the ordinary AR(1) least-squares estimator,
which the clean-statistics test of \code{tests/test\_em\_bp.py} confirms. The Gaussian
kernel's \code{m\_step} in \code{src/markov\_score/kernels.py} does not clamp
$|\corr|<1$: the update is the exact maximiser, a finite chain with $|\corr|\ge1$ is a
valid joint law, and what fails there is only the stationary interpretation
$\ivar=1-\corr^2$, which the code flags with a warning rather than silently biasing the
step.

\section{The mixture-innovation kernel}

\subsection{Definition and moments}

The paper's headline estimator keeps the linear autoregression and frees the
innovation law:
\begin{equation}
  \kernelth(a'\mid a)=\sum_{c=1}^{\ncomp}\pi_c\,\normal{a'-\corr a-\nu_c}{s_c^2},\qquad
  \params=(\corr,\pi_1,\dots,\pi_\ncomp,\nu_1,\dots,\nu_\ncomp,s_1^2,\dots,s_\ncomp^2),
  \label{eq:en-mixture}
\end{equation}
a $\ncomp$-component Gaussian mixture as the density of the innovation
$\varepsilon=a'-\corr a$. Heavy tails (Laplace, Student-$t$), bimodality and skew are all
approximable; the estimator never sees a Laplace density and must reproduce one from
noised data. The component means are unconstrained: recentering them after each
update to enforce $\sum_c\pi_c\nu_c=0$ is not a maximiser of $\Qfun$ and was found to
break monotone ascent, so the fitted innovation mean is reported as a diagnostic
instead (on zero-mean data it lands near zero by itself). The fitted innovation's
moments are read from the components in closed form,
\begin{equation}
\begin{aligned}
  m_1&=\sum_c\pi_c\nu_c,\qquad
  \ivar_{\mathrm{fit}}=\sum_c\pi_c\bigl(s_c^2+(\nu_c-m_1)^2\bigr),\\
  m_4&=\sum_c\pi_c\bigl(3s_c^4+6s_c^2(\nu_c-m_1)^2+(\nu_c-m_1)^4\bigr),\qquad
  \hat\kappa_4^{\mathrm{ex}}=\frac{m_4}{\ivar_{\mathrm{fit}}^2}-3,
\end{aligned}
  \label{eq:en-moments}
\end{equation}
using the Gaussian moments $\E(\nu+sZ-m_1)^4$ expanded in powers of $sZ$ (the odd
terms vanish). The last quantity is the fitted excess kurtosis, the scalar
``non-Gaussianity knob'' the paper compares with the truth ($+3$ for Laplace); the
formula is \code{innovation\_moments} in \code{src/markov\_score/kernels.py} and is
tested against draws from the same mixture.

\subsection{The complete-data problem has a second latent variable}

For the Gaussian kernel the M-step was a closed-form maximisation because
$\log\kernelth$ was a quadratic in $\params$. For \cref{eq:en-mixture} it is the
logarithm of a sum, and the maximiser of $\langle\Xis,\log\kernelth\rangle$ has no
closed form. The remedy is the classical one: the mixture has its own latent variable,
the component label, and given $\Xis$ the M-step is itself an EM problem over that
label. Define the \term{responsibilities} on the grid,
\begin{equation}
  r_c(k,l)=\frac{\pi_c\,\normal{\gridpt{k}-\corr\gridpt{l}-\nu_c}{s_c^2}}
  {\sum_{c'}\pi_{c'}\,\normal{\gridpt{k}-\corr\gridpt{l}-\nu_{c'}}{s_{c'}^2}},
  \label{eq:en-resp}
\end{equation}
the posterior probability of component $c$ given the transition $\gridpt{l}\to\gridpt{k}$.
A sweep of \term{conditional maximisations} (\cref{ch:maximum-likelihood-em}) then
alternates two blocks, each a closed-form maximiser of the inner Q-function with the
other block fixed, with $e_{kl}=\gridpt{k}-\corr\gridpt{l}$ the residual:
\begin{align}
  \text{mixture block:}\qquad
  \pi_c&=\frac{\langle\Xis,r_c\rangle}{W},\quad
  \nu_c=\frac{\langle\Xis,r_c\circ e\rangle}{\langle\Xis,r_c\rangle},\quad
  s_c^2=\frac{\langle\Xis,r_c\circ(e-\nu_c)^2\rangle}{\langle\Xis,r_c\rangle};
  \label{eq:en-mixblock}\\
  \text{autoregression block:}\qquad
  \corr&=\frac{\sum_{c,k,l}\Xis_{kl}\,r_c(k,l)\,\gridpt{l}\,(\gridpt{k}-\nu_c)/s_c^2}
              {\sum_{c,k,l}\Xis_{kl}\,r_c(k,l)\,\gridpt{l}^2/s_c^2},
  \label{eq:en-rhoblock}
\end{align}
the first being the usual Gaussian-mixture M-step with the weights $\Xis_{kl}r_c(k,l)$
and the second a weighted least-squares fit of the child on the parent with
component-specific offsets $\nu_c$ and precisions $1/s_c^2$ (differentiate
$\sum\Xis\,r_c[-(\gridpt{k}-\corr\gridpt{l}-\nu_c)^2/2s_c^2]$ in $\corr$). The
responsibilities are recomputed between the blocks. Each block increases the inner
Q-function, which by Jensen's inequality (the standard mixture-EM argument applied
with $\Xis$ as the ``data weights'') increases $\langle\Xis,\log\kernelth\rangle$; so the
sweep ascends $\Qfun$, and by the monotonicity theorem of
\cref{ch:maximum-likelihood-em} the outer EM's marginal likelihood cannot decrease.
This is the sense of the paper's ``a generalised EM whose ascent is monotone'' and of
its citation of Meng and Rubin.

The implementation, the mixture kernel's \code{m\_step} in
\code{src/markov\_score/kernels.py}, sweeps until the per-transition gain in
$\langle\Xis,\log\kernelth\rangle/W$ falls below $10^{-8}$ or sixteen sweeps are spent,
and records which happened on the returned kernel; a negative gain beyond round-off
raises an error, since a conditional maximiser cannot decrease its objective. The
docstring records why sixteen: with four sweeps the innovation \emph{shape} was
materially under-converged on short outer runs while $\corr$ barely moved, the same
asymmetry as the outer loop's (\cref{sec:en-asymmetry}), and sixteen removes the
interaction at a quarter of the cost of sixty-four. Because $\Xis$ is fixed during the
sweeps, each costs $O(\ncomp\ngrid^2)$ on a matrix that never grows with the dataset,
which is why the M-step's cost per iteration is independent of $\nseq$ and why, at the
operating grid, it is the M-step rather than the E-step that dominates a fit below a
few thousand sequences.

\subsection{Resolution}

A mixture component can be narrower than a grid cell: the variance floor in the code
is $10^{-6}$, a standard deviation of $10^{-3}$, against a spacing of $0.04$. A spike
narrower than the mesh can raise the quadrature likelihood without describing any
feature of the innovation law. The diagnostic \code{scale\_diagnostics} reports the
narrowest fitted standard deviation in units of $\gridh$ and a flag
\code{resolved} equal to $s_{\min}\ge2\gridh$; every reported cell carries it, the paper's
Appendix C says such cells are flagged rather than dropped, and its capacity table
marks them.

\section{The checkpoint ladder and validation selection}

\subsection{The ladder}

The protocol (\code{experiments/protocol.py}) fixes an iteration ladder
\[
  \{10,\,20,\,40,\,60,\,80,\,120,\,160,\,220,\,300,\,400,\,600,\,800,\,1200\}.
\]
\code{fit\_em} runs EM to the top rung and snapshots the kernel at each rung. The snapshot semantics are precise: at
iteration $k$ the kernel about to be scored is saved \emph{before} its E-step, so the
saved kernel at rung $k$ is the one whose likelihood is the $k$-th trace entry; the
final iterate is always available under its own index, so a caller asking for a rung
past an early convergence break still receives the best kernel rather than an error.
The likelihood tolerance ($10^{-9}$ nats per transition, signed) can end a run before
the top; at the paper's sizes it rarely does, and the run's \code{stop\_reason} records
which happened.

\subsection{Selection and the \code{at\_cap} flag}

Which rung is reported is chosen on a \term{validation bundle}: sequences drawn from a
random stream disjoint from the training stream, noised, and scored by the relative
score error on the bulk sites (the four sites at each end trimmed) against the
reference posterior means computed with the \emph{true} kernel. The rung with the
smallest validation error is selected and then scored, once, on a separate test
bundle. The selected rung is compared with the top of the ladder, and the boolean
\code{at\_cap} records whether they coincide (\code{best == max(saved)} in
\code{experiments/e02\_capacity.py} and \code{experiments/e06\_learning\_through\_noise.py}).
A cell with \code{at\_cap} true is one whose validation curve was still falling at the
last rung, so its error is an upper bound on what a longer ladder would give; the
paper reports the fraction of such cells (as the macro that its Section 4 calls the
at-cap rate) and refuses to read a difference between two arms as final when either
is at cap. E20 shows that the oracle in this rule is dispensable: selecting the rung
by the held-out one-view log-evidence per transition, which needs neither the clean
chain nor the true kernel, picks the same rung in most seeds and the same test error
within noise.

\section{The clean-data limit}
\label{sec:en-clean}

As $t\to0$ the unary factors sharpen to delta functions at the observations, the
posterior collapses onto the clean chain, the pairwise beliefs become products of
deltas, and $\Xis$ degenerates to a histogram of the observed transitions. The code
makes this limit an explicit input: \code{clean\_statistics} in
\code{src/markov\_score/em.py} places each observed pair $(a_{i-1},a_i)$ on the grid by
bilinear interpolation (a value $a$ between $\gridpt{k}$ and $\gridpt{k+1}$ is split into
weights $1-\theta$ and $\theta$ with $\theta=(a-\gridpt{k})/\gridh$, which preserves the
first moment exactly), and accumulates unit mass per transition. No quadrature weights
enter: the clean $\Xis$ is a counting measure, whereas the BP $\Xis$ is belief mass that
has already absorbed the weights; both sum to the number of transitions, which is the
only property the M-steps rely on. \code{fit\_clean} then applies the same M-step
code, so the clean-data maximum-likelihood fit and the noisy-data EM differ only in
how $\Xis$ is built. The test \code{test\_small\_noise\_em\_approaches\_clean\_mle} pins
the continuity between the two. E06 includes a clean arm as the ``optimiser's own
floor''; for one code path it is implemented as the channel at a tiny level, which
the paper's Figure 3 shows as the dotted reference.

\section{How the experiments use this}

\paragraph{E02a (capacity).} Laplace chain at the frozen configuration; for each seed
and size $\nseq$ the same hashed bundles; mixture kernels with
$\ncomp\in\{1,2,4,8,16\}$ from the same initialisation stream ($\corr$ initialised at
$0.3$, components spread over the innovation range), one view per sequence split over
the twelve levels, the ladder above, the rung selected on validation identically for
every capacity, then the test error in every region, the fitted $\corr$, innovation
variance and excess kurtosis, the resolution flag and \code{at\_cap}. The $\ncomp=1$ arm
must land on the LMMSE floor by the corollaries of \cref{ch:gaussian-closure-lmmse};
the paper's capacity table (Section 5) reports that it does, and that the mixture goes
far below it. The
earlier reading that ``capacity does not help'' came from a sweep in which only the
single-Gaussian arm stopped at its optimum (the tolerance fired for it) while the richer
arms ran to a fixed cap past theirs: the stopping rule, not the model, produced the
verdict (the paper's Appendix F).

\paragraph{E06 (learning through noise).} The same estimator with $\ncomp=8$ and
$\nseq=512$, but with every training sequence noised at \emph{one} level $t$, over a
ladder of levels and sixteen replicates, the validation bundle at the same level, the
test bundle on the twelve-level schedule. Recorded per cell: the selected rung,
\code{at\_cap}, the stop reason, $\hat\corr$, the fitted innovation variance and excess
kurtosis, the Hellinger distance of the fitted kernel to the truth, the resolution
flag and the schedule-pooled test error. The paper's Figure 3 and Section 4 read these
against the Cram\'er--Rao reference $1/\sqrt{\nseq I_t}$ level by level: unimodal and at
the reference at low levels; a split into a shrunk cluster and Gaussian-collapsed fits
at an intermediate level; collapse where the reference crosses the coordinate's range;
$\hat\corr$ recovered until past the crossover. At the two highest levels the selected
rung is the top of the ladder in most cells, which is the budget question.

\paragraph{E06c.} The control that answers it: the same cells with the ladder
extended sixfold. Its result is what decides whether the excess of $\hat\corr$'s error
over the reference at those levels is information or optimisation; the paper leaves
the question open pending it, and \cref{sec:en-asymmetry} says why the budget binds
exactly there.

\section{The convergence asymmetry}
\label{sec:en-asymmetry}

The paper reports, as a fact about this model, that the correlation coordinate
settles in tens of EM updates and the innovation shape in hundreds. Two mechanisms
produce it, and both can be read off the information at each Hermite degree.

\emph{The outer loop.} By the rate formula of \cref{ch:maximum-likelihood-em}, EM
contracts the error along a coordinate by the factor
$1-I_{\mathrm{obs}}/I_{\mathrm{com}}$ per iteration, where $I_{\mathrm{obs}}$ is the
information the noised data carry about the coordinate and $I_{\mathrm{com}}$ the
information the clean chains would carry. The complete information is of the same
order for both coordinates: clean transitions determine $\corr$ and the innovation
shape at the same rate in $\nseq$. The observed information differs by the Hermite
degree at which each coordinate first appears: by the paper's decay theorem, per sequence
it is $\tfrac12\Frob{\partial_\corr\Sig_0}^2e^{-4t}$ for the correlation (degree two) and
$\tfrac1{24}\Frob{\partial_\theta\kfour}^2e^{-8t}$ for a shape coordinate (degree four), up
to corrections. So the missing fraction is closer to one for the shape at every level,
by roughly a factor $e^{-4t}$ in the retained fraction ($0.37$ at $t=0.25$, $0.14$ at
$t=0.5$, $0.018$ at $t=1$), and the number of iterations the shape needs grows
correspondingly. In the clean limit both informations are complete and the asymmetry
disappears, which the clean arm of E06 shows.

\emph{The inner loop.} For the mixture kernel there is a second latent variable, the
component label, and the shape parameters are learned through it: the M-step's inner
sweeps are themselves an EM whose rate is governed by how well the responsibilities
\cref{eq:en-resp} separate the components. A nearly Gaussian innovation gives nearly
equal responsibilities, a missing fraction near one, and slow inner convergence,
which the code's sixteen-sweep budget was introduced to absorb; $\corr$ is not learned
through the label at all (its block is a single least-squares solve) and does not
suffer this. The two mechanisms compound, and both act on the same coordinate.

The consequence for the one-view protocol is the paper's arithmetic: of the twelve
levels, spread over $[0.05,3]$, only the lowest few carry appreciable information about
the shape, since $e^{-8t}$ has fallen by a factor of several hundred by $t=0.7$, so
for the shape the twelve-level schedule is effectively a four-level one, and the
estimator's shape is learned from a quarter of the data at a rate set by the smallest
of its information ratios. That is why the ladder is long, why the shape is what the
inner budget affects, and why the validation-selected rung sits at the cap precisely
at the levels where the information is smallest.

\section{A worked example: the Gaussian M-step from a small $\Xis$}

Return to the three-site, three-state chain of \cref{ch:factor-graphs-bp}, grid
$\{-1,0,1\}$. Its two pairwise posteriors summed over the two edges give, for that one
sequence,
\[
  \Xis=\begin{pmatrix}0.0802&0.1585&0.0665\\0.0661&0.7356&0.2385\\0.0251&0.2638&0.3656\end{pmatrix},
  \qquad W=\sum_{kl}\Xis_{kl}=2,
\]
rows indexed by the child $\gridpt{k}$, columns by the parent $\gridpt{l}$ (the first edge
matrix was written out in that chapter; the second is obtained the same way). The
belief-weighted moments are
\begin{align*}
  S_{00}&=\sum_{kl}\Xis_{kl}\gridpt{l}^2=(0.0802+0.0661+0.0251)+(0.0665+0.2385+0.3656)=0.8421,\\
  S_{11}&=\sum_{kl}\Xis_{kl}\gridpt{k}^2=(0.0802+0.1585+0.0665)+(0.0251+0.2638+0.3656)=0.9598,\\
  S_{01}&=\sum_{kl}\Xis_{kl}\gridpt{k}\gridpt{l}=0.0802-0.0665-0.0251+0.3656=0.3542,
\end{align*}
where $S_{01}$ collects the four corner entries with the sign of $\gridpt{k}\gridpt{l}$. Then
\cref{eq:en-yw} gives $\corr=0.3542/0.8421=0.4206$ and
$\ivar=(0.9598-0.4206\cdot0.3542)/2=0.4054$. Evaluating $\Qfun(\corr,\ivar)$ of
\cref{eq:en-Qgauss} at this point gives $-1.9350$, and at the perturbed points
$(0.4706,\,0.4054)$, $(0.4206,\,0.4459)$ and $(0.3706,\,0.3649)$ it gives $-1.9376$,
$-1.9394$ and $-1.9436$: the closed form is the maximiser, as the test of the Gaussian
M-step in \code{tests/test\_em\_bp.py} checks by random perturbation. A two-component
mixture started at $\pi=(\tfrac12,\tfrac12)$, $\nu=(-0.3,0.3)$, $s^2=(0.5,0.5)$,
$\corr=0.2$ has $\langle\Xis,\log\kernelth\rangle/W=-1.0155$; one mixture block
\cref{eq:en-mixblock} moves it to $-0.9700$ and the following autoregression block
\cref{eq:en-rhoblock} to $-0.9592$, with $\corr=0.328$ after one sweep: each block
increases the objective, as the conditional-maximisation argument requires. (The
example has a single sequence and a three-point grid, so the numbers illustrate the
algebra, not the statistics.)

\where{The paper, Section 4 (the marginal likelihood, Fisher's identity and equation
\emph{fisher}, the Yule--Walker M-step, the mixture kernel and its
conditional-maximisation sweeps, the one-view protocol and validation selection, the
convergence asymmetry and the twelve-level arithmetic), Appendix C (the ladder and its
selection rule, the resolution flag), Appendix E (E19's moment estimators as the
competitor, E20's oracle-free selection), Appendix F (the stopping-rule artefact of the
earlier capacity reading). Code: the E-step, its multi-level sum, the clean-data
statistics, the two fitting loops and the Q-function helpers are all in
\code{src/\allowbreak markov\_score/\allowbreak em.py}; the three kernels, with the mixture's
\code{innovation\_moments} and \code{scale\_diagnostics}, are in
\code{src/\allowbreak markov\_score/\allowbreak kernels.py}; the ladder \code{em\_checkpoints} is in
\code{experiments/\allowbreak protocol.py}; the drivers are \code{experiments/\allowbreak e02\_capacity.py},
\code{experiments/\allowbreak e06\_learning\_through\_noise.py} and
\code{experiments/\allowbreak e20\_evidence\_selection.py}. Tests: \code{tests/\allowbreak test\_em\_bp.py},
\code{tests/\allowbreak test\_fisher\_bruteforce.py} and
\code{tests/\allowbreak test\_gradient\_trace\_alignment.py}.}

\paragraph{References.}
A.~P.~Dempster, N.~M.~Laird and D.~B.~Rubin (1977), ``Maximum likelihood from incomplete
data via the EM algorithm'', \emph{Journal of the Royal Statistical Society B} 39(1).
X.-L.~Meng and D.~B.~Rubin (1993), ``Maximum likelihood estimation via the ECM algorithm:
a general framework'', \emph{Biometrika} 80(2).
O.~Capp\'e, E.~Moulines and T.~Ryd\'en (2005), \emph{Inference in Hidden Markov Models},
Springer, Chapters 10 and 11 (Fisher's identity, EM for HMMs).
R.~H.~Shumway and D.~S.~Stoffer (1982), ``An approach to time series smoothing and
forecasting using the EM algorithm'', \emph{Journal of Time Series Analysis} 3(4).
R.~Douc, E.~Moulines, J.~Olsson and R.~van Handel (2011), ``Consistency of the maximum
likelihood estimator for general hidden Markov models'', \emph{Annals of Statistics} 39(1).
